\begin{prob}[(3 points)]

    Consider the following problem: given two \textbf{sorted} lists, \python{A}
    and \python{B}, each containing $n$ numbers, determine if the two lists
    share an element in common. That is, determine if there is a number $x$
    which appears in both lists.

    \begin{subprobset}

        \begin{subprob}
            What is a \textbf{tight} theoretical lower bound for this problem?
            Again, you may assume that the lists are both sorted.
            State your answer in asymptotic notation as a function of $n$.

            \inlineresponsebox[.75in]{$\Theta(n)$}
        \end{subprob}

        \begin{subprob}(2 points)
            Briefly (in 50 words or less) describe an algorithm for solving
            this problem \textbf{and state its worst case time complexity}. To
            earn full credit, you should give the optimal algorithm, but
            partial credit will be awarded for sub-optimal solutions. You may
            write pseudocode if you like, but an explanation in words is also
            sufficient. Make sure that you only state one algorithm!

            \begin{responsebox}{2in}
                Here's an approach that takes $\Theta(n)$; it's similar to what we did in lecture
                for the movie problem.

                Start with \python{i = j = 0}. If \python{A[i] == A[j]}, return \python{True}, since
                you've found an element in common. If \python{A[i] > A[j]}, increment \python{j}, otherwise
                increment \python{i}. Takes $\Theta(n)$ time.

                The brute force approach also earns credit: Loop over all pairs, checking if 
                \python{A[i] == B[j]} for each pair. This takes $\Theta(n^2)$ time.

            \end{responsebox}

        \end{subprob}

    \end{subprobset}

\end{prob}
